% !TeX root = ../papers/nc-p3.tex
\section{五つの四次中心切断の像と重み付き完全交叉}
\label{sec:quartic-center-image-atlas}
\begin{theorem}\label{thm:quartic-center-image-atlas}
前節の十五次元閉包の正則整域窓で、五次元の $Z_1$ が定める有理写像の像閉包を
$Y\subset\mathbb P^4$ とする。次の四型で全てを覆う。
\begin{enumerate}
\item $Y$ は四次三次元超曲面で、$Z$ はその通常の座標環である。
\item $Y$ は三次元二次曲面で、$Z=k[x_0,\ldots,x_4,w]/(q_2,F_4)$、
$\deg x_i=1,\deg w=2$ である。四次中心系に基点がある場合、
$F_4=w l_2(x)+f_4(x)$、$l_2$ は $k[x]/q_2$ で非零となり、写像は双有理である。
\item $Y$ は二つの二次式の正規な完全交叉曲面である。
$Z=k[x_0,\ldots,x_4,u,w]/(q_1(x),q_2(x),F_4(x,u,w))$、
$u,w$ は次数2、一般の $u$ を八次切断に取って $F_4$ は $w$ に関してモニックな二次式となる。
一般八次曲面の標準写像は $Y$ への有限二重被覆である。
\item $Y$ は次数3の非退化曲面（有理正規スクロールまたはその錐）である。
一般八次曲面の標準系に固定曲線はなく、基点は滑らかな部分の異なる単純二点である。
一般標準曲線は種数9の超楕円曲線で、この二点はいずれもそのWeierstrass点である。
\end{enumerate}
ここで各型の中心を実現する窓の存在を任意の係数で主張してはいない。
\end{theorem}
\begin{proof}[標準像の次元と次数]
$S$ を幾何学的一般八次曲面とする。$H^0(X,H)\to H^0(S,K_S)$ は同型である。
切断列と深さ4、$K_S=H|_S$ による。前節の横断 $A_1$ 点では $|K_S|$ に基点がない。
実際定理\ref{thm:index-four-quartic-basepoints}で非Azumayaの基点は必ず安定化点で、$S$ はそれを避ける。
$S$ の最小解消の標準因子 $K$ は $K_S$ の引き戻しで、ネフかつ $K^2=8$ である。
もし標準系の像が曲線なら、そのStein分解の一般繊維の最小解消上の移動曲線を $C$ とする。
非退化曲線の次数が少なくとも4なので $K$ の一般因子には少なくとも四つのそのような曲線が入る。
$KC>0$, $C^2\ge0$ であり、もし $KC=1$ ならHodge指数から $C^2=0$、
随伴の偶奇に反する。従って $KC\ge2$、等号と四繊維だけの場合しかない。
残る固定部の $K$ との積は零なので例外曲線上だけであるが、それらの場所は基点でなかった。
よって $K\equiv4C$ となり、$2=KC=4C^2$ に矛盾する。像が一点でないことは五切断の独立性から従う。

三次元 $X$ が曲面 $Y$ を支配する場合、一般 $S$ は今や $Y$ を生成有限に支配する。
その次数 $d$ は少なくとも2である。実際 $X$ の滑らかな固有解消には大域二形式がない。
$X$ は有理特異点を持つ。これは有理特異点を持つ中心錐のエタールな次数自明化と、
切片の有限群商での有理特異点の保存（標数零の商の定理）から従う。
$H^2(X,\mathcal O_X)=0$ は中心の深さ4によるから、滑らかな解消でのHodge対称性を使えばよい。
以上は標数零で有限型のデータを複素数に移して検査できる。$d=1$ なら
$Y$ の滑らかなモデルに $p_g(S)=5$ 個の二形式があり、$X$ の解消への有理写像による引き戻しに反する。
一方 $d\deg Y\le K^2=8$。基点を解消した自由な移動部の平方が左辺で、固定部を除いても増えない
（解消上で $K$ の引き戻しと移動部はともにネフ）。非退化曲面は次数3以上なので $d=2$、次数は3または4である。

$Y$ が三次元なら非退化な超曲面である。次数2以上の超曲面で二次式を持たなければ、
$\dim Z_2=15$ により $Z_2=\operatorname{Sym}^2 Z_1$ となる。八次系は基点無しなので
四次系も基点無しで、既出の四中心パラメータによる分類から第一型となる。
$Y$ が二次曲面で基点無しの場合も同分類で第二型となる。基点ありとする。
有理写像を解消し、$A$ を $H$ の引き戻し、$M$ を $Y$ の超平面の引き戻しとする。
$2(A-M)$ は有効であり、両者はネフかつ巨大である。$M^3=2\deg(X/Y)\le A^3=4$。
等号なら混合積 $A^2M,AM^2$ も全て等しく、
ネフかつ巨大な因子のKhovanskii--Teissier不等式の等号条件で $A,M$ は数値的に等しい。
有効な差は零となり、解消上の二乗切断系に固定部がなくなる。
基点があればその基底イデアルの可逆化には非零の有効因子が出るので矛盾する。
従って次数は1。$Z_2$ に一つの補元 $w$ を取り、$D=k[Z_1,w]$ とする。
八次パラメータを含むから $Z$ は $D$ 上有限で、双有理性と次数1の元の存在により分数体は同じである。
$A_0=(k[x]/q_2)[w]$ のHilbert級数は $(1-t)^{-5}$ である。
$D\subset Z$ と $H_Z=(1-t^4)/(1-t)^5$ より $A_0\to D$ の核には非零四次式 $F_4$ がある。
$A_0/(F_4)$ と $Z$ の級数は一致し、$D$ の最高次の重複度も有限双有理性で同じである。
従ってこの完全交叉は素で $D$ となる。実際他の極小成分も生成点重複も重複度の加法性で許されず、
埋込み素因子もない。次元別の一致で $D=Z$。基点では全 $x_i=0$ かつ $w\ne0$ なので
$F_4$ に $w^2$ 項がなく、$Y$ をさらに切る関係はないから $w$ の係数は非零である。
\end{proof}

\begin{proof}[二つの曲面像]
次数4の場合、解消上の $K$ の引き戻しと移動部の平方がともに8である。
先と同じ二次元のHodge指数の等号条件で差は零、標準系は自由である。$K_S$ は豊富なので写像は有限である。
$Y^\nu$ を正規化、$L$ を超平面の引き戻しとする。$h^0(L)=5$（$S$ へ引き戻せる）、
$q(Y^\nu)=0$（$S$ からの有限traceで分裂）である。一般 $C\in|L|$ は滑らかな整曲線で、
$L|_C$ の全切断は四次元である。この次数4の線型系は $C$ を $\mathbb P^3$ の非退化四次曲線へ双有理に写す。
Castelnuovoの種数上界は1で、種数0なら完全四次系の次元に反する。よって種数1で
$h^0(2L)\le5+8=13$。従って $Y$ は独立な二二次式に入り、非退化整曲面で次数4ゆえ
その完全交叉である（共通超平面因子は $Y$ の非退化性から生じない）。一般の超平面曲線は
完全四次系による楕円曲線の埋込みで滑らかであるから、完全交叉の $Y$ は余次元1で正則、従って正規である。

$P$ をこの二二次式の標準次数付き座標環とする。$b\in Z_2$ を一般に取ると $T=Z/bZ$
は $S$ の標準環である。実際全ての次数で切断列が全射で、一般曲面には次数安定化がない。
$T$ は $P$ 上有限で分数体次数2、双方とも正規Cohen--Macaulayである。
traceで $T=P\oplus N$ と分ける。$\omega_P=P(-1)$ と $\omega_T=T(1)$ は
完全交叉とGorenstein切断から従う。有限双対の $\operatorname{Hom}_P(T,\omega_P)=\omega_T$ において、
$P$ 上の二次被覆の対合を考える。$\omega_T$ の斉次自由生成元への指標は $+1$ または $-1$ である。
$+1$ なら不変部分は $P(1)$ となって $\omega_P=P(-1)$ に一致しない。
よって $-1$ で、不変部分は $N(1)=P(-1)$、すなわち $N=P(-2)$。
従って $T=P[w]/(w^2-f_4)$ である。$w$ を $Z$ に持ち上げ、$u=b$ とする。
次数ごとに $u$ で割り戻せば $Z$ は $x_i,u,w$ で生成され、二つの二次式は既に $Z$ 内で零である。
$w^2-f_4$ は $u$ 倍の二次式なので、これを引いてモニック関係を得る。
得た完全交叉の級数は $Z$ と一致し、第三型となる。

次数3の場合は最小次数曲面の分類で上記のスクロール型となる。
$S$ の最小解消上で $K$ の固定曲線を合わせたものを $G$ とする。
標準写像をさらに解消した移動部 $M$ は $M^2=6$ なので $KM\ge\sqrt{48}$、従って $KM\ge7$。
$G\ne0$ なら例外固定曲線はないので $KG=1$、Hodge指数と随伴の偶奇で $G^2\le-1$ である。
すると $(K-G)^2\le5$ で $M^2=6$ に反する。よって固定曲線はなく、基点重複度の平方和は2。
それらは最小解消の例外曲線から離れ、一般標準曲線は単純基点でも滑らかである
（一般切断の多重度が1）。一般曲線は整で、種数9、スクロールの一般超平面である有理正規三次曲線への
二重被覆を持つ。この超楕円鉛筆を $J$ とする。標準系のこの曲線への制限は
$q=0$ から $K_S|_C$ の完全系で、次数8、次元3、移動部は $3J$、固定部 $P+Q$ である。
随伴から $2K_S|_C=K_C=8J$、従って $2P+2Q\sim2J$。
$|2J|$ は二繊維の和だけなので、$P+Q$ 自体が一繊維でなければ $P,Q$ は異なる二分岐点である。
一繊維なら固定部分に動く鉛筆が含まれて矛盾する。従って無限近接基点もなく第四型を得る。
\end{proof}
